\documentclass[border=10pt,tikz]{standalone}
\usepackage{fontspec}
\setmainfont{Apple SD Gothic Neo}
\setsansfont{Apple SD Gothic Neo}
\usepackage{tikz}
\usetikzlibrary{positioning,arrows.meta,shapes,fit,backgrounds,calc}
\begin{document}
% _design.tex — shared TikZ design system for all blog diagrams
% Inputs nothing on its own — meant to be \input{} from each diagram .tex
% AFTER \begin{document}, BEFORE \begin{tikzpicture}.
%
% Provides a unified visual language across GoF / DSA / EMC++ / Beautiful series:
%   - Colors: muted, accessible, color-coded by role
%   - Node styles: consistent sizing, rounded, subtle borders
%   - Edge styles: UML-style inheritance, aggregation, plus dashed alternates
%
% Usage:
%   \documentclass[border=10pt,tikz]{standalone}
%   \usepackage{tikz}
%   \usetikzlibrary{positioning,arrows.meta,shapes,fit,backgrounds}
%   \begin{document}
%   \input{../_design.tex}     % adjust relative path
%   \begin{tikzpicture}[blog]  % apply 'blog' style
%   ...

\definecolor{absfill}{RGB}{220,232,247}   % cool blue
\definecolor{absbord}{RGB}{72,118,182}
\definecolor{confill}{RGB}{249,222,222}   % warm red
\definecolor{conbord}{RGB}{197,93,93}
\definecolor{impfill}{RGB}{223,238,222}   % soft green
\definecolor{impbord}{RGB}{99,154,93}
\definecolor{clifill}{RGB}{251,239,205}   % pale yellow
\definecolor{clibord}{RGB}{200,167,72}
\definecolor{hifill} {RGB}{253,224,200}   % accent orange
\definecolor{hibord} {RGB}{210,135,68}
\definecolor{nodefill}{RGB}{250,250,250}
\definecolor{nodebord}{RGB}{120,120,120}
\definecolor{edgec}  {RGB}{80,80,80}
\definecolor{edgesc} {RGB}{145,145,145}
\definecolor{groupbord}{RGB}{200,200,200}

\tikzset{
  blog/.style={
    font=\sffamily\small,
    every node/.style={font=\sffamily\small}
  },
  % --- node roles ---
  cls/.style={
    rectangle, rounded corners=2pt,
    draw=nodebord, line width=0.4pt,
    minimum height=8mm, minimum width=28mm,
    inner sep=3pt, fill=nodefill, align=center
  },
  abs/.style={cls, draw=absbord, fill=absfill, font=\sffamily\itshape\small},
  con/.style={cls, draw=conbord, fill=confill},
  imp/.style={cls, draw=impbord, fill=impfill},
  cli/.style={cls, draw=clibord, fill=clifill},
  hi/.style={cls, draw=hibord, fill=hifill, font=\sffamily\bfseries\small},
  % --- circle (for trees / graph nodes) ---
  vx/.style={
    circle, draw=absbord, line width=0.4pt,
    minimum size=8mm, fill=absfill, inner sep=0pt
  },
  vxc/.style={vx, draw=conbord, fill=confill},
  vxh/.style={vx, draw=hibord, fill=hifill, font=\sffamily\bfseries},
  % --- decision diamond ---
  q/.style={
    diamond, draw=clibord, line width=0.4pt, fill=clifill,
    inner sep=2pt, aspect=2.4, align=center
  },
  % --- edges ---
  inh/.style={                                  % UML inheritance
    -{Triangle[length=3mm,width=3mm,open]},
    draw=edgec, line width=0.5pt
  },
  agg/.style={                                  % UML aggregation
    {Diamond[length=2.5mm,width=2.5mm,open]}-{Stealth[length=2.2mm]},
    draw=edgec, line width=0.45pt
  },
  uses/.style={                                 % directed reference
    -{Stealth[length=2.2mm]}, draw=edgec, line width=0.5pt
  },
  arr/.style={                                  % plain arrow (graphs/trees)
    -{Stealth[length=2mm]}, draw=edgec, line width=0.45pt
  },
  edge/.style={                                 % undirected (trees)
    draw=edgec, line width=0.45pt
  },
  alt/.style={                                  % dashed: similar/alt
    -{Stealth[length=2mm]}, dashed, draw=edgesc, line width=0.45pt
  },
  ret/.style={                                  % dashed creates/returns
    -{Stealth[length=2mm]}, dashed, draw=edgesc, line width=0.45pt
  },
  thr/.style={                                  % threading link (red dashed)
    ->, dashed, draw=conbord, line width=0.5pt
  },
  % --- group / region ---
  group/.style={
    rectangle, rounded corners=4pt,
    draw=groupbord, line width=0.4pt,
    inner sep=10pt, fill=black!2
  },
  glabel/.style={font=\sffamily\bfseries\small, anchor=north west, inner sep=2pt}
}

% Korean font convenience macro — included by every Hangul diagram
\providecommand{\KOR}{}

\begin{tikzpicture}[blog,
  bx/.style={cls, minimum width=52mm, text width=48mm, minimum height=22mm, align=center, font=\sffamily\footnotesize, inner sep=4pt},
  rom/.style={bx, draw=absbord, fill=absfill},
  ub/.style={bx, draw=impbord, fill=impfill},
  mem/.style={bx, draw=hibord, fill=hifill},
  lay/.style={bx, draw=clibord, fill=clifill},
  hdr/.style={font=\sffamily\bfseries\small, text=inkc, anchor=west},
  sub/.style={font=\sffamily\footnotesize, text=mutec, anchor=west}
]

% ── Stage 1: ROM → first firmware ──
\node[hdr] at (-2.8, 7.6) {① ROM에서 다음 firmware로};
\node[sub] at (-2.8, 7.05) {근거: SoC reference manual · 보드 문서};

\node[rom] (sel) at (0, 5.4) {%
  \textbf{boot source 선택}\\[2pt]
  reset 설정 · straps ·\\[1pt] 비휘발성 설정};
\node[lay] (img1) at (6.6, 5.4) {%
  \textbf{첫 firmware stage 이미지}\\[2pt]
  ROM 전용 header · 고정 위치 ·\\[1pt] packaging format (SoC별)};
\node[mem] (fw) at (13.2, 5.4) {%
  \textbf{다음 firmware 실행}\\[2pt]
  \scriptsize (이후 DRAM 초기화 등, 9장)};

\draw[flow] (sel) -- node[font=\sffamily\scriptsize, text=mutec, above] {읽음} (img1);
\draw[flow] (img1) -- (fw);

% ── Stage 2: U-Boot → OS files ──
\node[hdr] at (-2.8, 2.6) {② U-Boot에서 OS 파일로};
\node[sub] at (-2.8, 2.05) {근거: board defconfig · DTS · U-Boot 버전별 command 문서};

\node[ub] (drv) at (0, 0.2) {%
  \textbf{저장 장치 driver · command}\\[2pt]
  \scriptsize eMMC · SD → \texttt{mmc}\\[1pt]
  \scriptsize SATA (SCSI 경로일 때) → \texttt{scsi}\\[1pt]
  \scriptsize raw NAND → \texttt{nand} · SPI NOR → \texttt{sf}};
\node[lay] (img2) at (6.6, 0.2) {%
  \textbf{저장 레이아웃}\\[2pt]
  raw offset · GPT/MBR 파티션 ·\\[1pt] filesystem · platform container};
\node[mem] (ram) at (13.2, 0.2) {%
  \textbf{RAM에 적재}\\[2pt]
  kernel · DTB · initramfs\\[2pt]
  \scriptsize 예: \texttt{load mmc <dev>:<part> ...}};

\draw[flow] (drv) -- node[font=\sffamily\scriptsize, text=mutec, above] {읽기} (img2);
\draw[flow] (img2) -- (ram);


\node[caption, text width=160mm] at (6.6, -1.9) {ROM이 읽는 위치와 U-Boot가 OS 파일을 읽는 위치는 따로 문서화합니다. 주소 · offset · Fuse 값은 SoC와 보드별 자료로 확인합니다.};

\end{tikzpicture}
\end{document}
